
\renewcommand{\arraystretch}{1.5} % Default value: 1

%%% Write here
%%%%%%%%%%%%%%%%%%%%%%%%%%%%%%%%%%%%%
\section{Introduction}
The Faraday effect, also known as Faraday rotation, is a magneto-optic phenomenon that demonstrates the fundamental interaction between light and a magnetic field. Discovered by Michael Faraday, this effect reveals that isotropic media become optically active when a magnetic field is applied parallel to the propagation direction of the light. This pivotal observation served as one of the earliest indications of the deep, underlying connection between optics and electromagnetism. Specifically, introducing a magnetic field in the direction of a light beam's propagation through a transparent medium induces a phenomenon known as circular birefringence.

As linearly polarized light travels through the optical medium under these conditions, its plane of polarization undergoes a measurable rotation. The angle of this rotation is a linear function of both the mean magnetic flux density and the length of the optical medium. The factor of proportionality governing this relationship is known as Verdet's constant. This specific experiment focuses on measuring and interpreting this magnetically induced circular birefringence.
%%%%%%%%%%%%%%%%%%%%%%%%%%%%%%%%%%%%%
\section{Objective}
The objective of this experiment is to measure and interpret the magnetically induced circular birefringence known as the Faraday effect.
%%%%%%%%%%%%%%%%%%%%%%%%%%%%%%%%%%%%%
\section{Equipment Required}
Here is the list of required equipment for the Faraday effect setup:
\begin{enumerate}
    \item \textbf{Light Source:} Red diode lasers with wavelengths of 650 nm, secured on a laser mount.
    \item \textbf{Polarizing Filters:} A polarizer rotator to polarize the incident light and a precision polarizer mount acting as an analyzer.
    \item \textbf{Magnetic Field Generator:} An electromagnet formed by a long solenoid with 2508 turns of wire, connected to a magnet power supply. Number of turns per unit length is 167.2 /cm.
    \item \textbf{Optical Medium:} A transparent flint glass rod, which serves as the sample placed inside the solenoid.
    \item \textbf{Detection System:} A photodiode detector used to measure the intensity of the light beam after it passes through the analyzer.
    \item \textbf{Mounting Apparatus:} A linear optical rail equipped with translating mounts. This track facilitates the secure attachment and precise spatial alignment of all optical components, ensuring the beam remains centered through the sample and upon the active area of the detector.
\end{enumerate}
%%%%%%%%%%%%%%%%%%%%%%%%%%%%%%%%%%%%%
\section{Theory}
%%%%%%%%%%%%%%%%%%%%%%%%%%%%%%%%%%%%%
\section{Procedure}
%%%%%%%%%%%%%%%%%%%%%%%%%%%%%%%%%%%%%
\section{Data analysis and calculation} 
%%%%%%%%%%%%%%%%%%%%%%%%%%%%%%%%%%%%%
\section{Conclusion}
%%%%%%%%%%%%%%%%%%%%%%%%%%%%%%%%%%%%%